\begin{solapartat}
\punts{0,4} Al primer cilindre càrrega negativa per tal de repel·lir l'electró i al segon cilindre càrrega positiva per tal d'atreure'l.

\punts{0,5} Les línies de camp surten de les càrregues positives i van a parar a les càrregues negatives:
\figura[0.4]{sol_fig1.png}

\punts{0,25} $|\vec{E}| = \dfrac{\Delta V}{d} \Rightarrow d = \dfrac{\Delta V}{|\vec{E}|} = \dfrac{250.000}{8,00\times 10^{6}} = 0,0312\ \mathrm{m} = 3,12\times 10^{-2}\ \mathrm{m}$ \quad \punts{0,1}
\end{solapartat}

\begin{solapartat}
\punts{0,6} $E_c = n\cdot e\cdot\Delta V = 1,00\times 10^{6}\ \mathrm{eV} = 1,602\times 10^{-13}\ \mathrm{J}$
\[ n = \frac{E_c}{e\cdot\Delta V} = \frac{1,602\times 10^{-13}}{1,602\times 10^{-19}\cdot 250.000} = 4 \]
O de la definició de eV, directament:
\[ n = \frac{E_c}{e\cdot\Delta V} = \frac{1,00\times 10^{6}\ \mathrm{V}}{250.000} = 4 \]

\punts{0,4} $v = \sqrt{2\,\dfrac{E_c}{m_e}} = 5,93\times 10^{8}\ \mathrm{m/s}$

\punts{0,25} Surt un resultat impossible $v > c$, la velocitat de l'electró no pot superar la velocitat de la llum. Caldria aplicar la correcció relativista per calcular la velocitat correcte.
\end{solapartat}
